% !TeX root = ../../Book2.tex
% 纲要标题：### D.4 von Neumann 代数与双交换子

% 纲要要点（供目录核对）：
% * 弱算子拓扑、强算子拓扑与范数拓扑；含恒等算子的 *-子代数的闭包
% * 双交换子定理的含幺自伴版本：\(A''=\overline A^{\mathrm{SOT}}=\overline A^{\mathrm{WOT}}\)
% * 中心、因子、投影与迹；以投影等价、有限性及极小性介绍 I、II、III 型因子及 II_1、II_∞ 的区别，计算 \(B(H)\) 的投影模型；不展开完整分类与 Tomita–Takesaki 模理论
% #### D.4.1 弱算子、强算子与范数拓扑
% #### D.4.2 双交换子定理
% #### D.4.3 中心、因子、投影与迹
% #### D.4.4 投影等价与因子的类型

\section{von Neumann 代数与双交换子}
\label{b2:sec:D04}

当 Hilbert 空间 \(H\) 有限维时，\(B(H)\) 上的范数、强算子与弱算子拓扑一致；无限维时则必须区分怎样取极限。双交换子定理的内容，正是把一种算子拓扑的闭包与代数交换关系联系起来。进入弱闭代数以后，还可以通过其中投影的比较，区分有限维矩阵没有表现出的结构。本节始终令 \(H\ne0\) 为复 Hilbert 空间，所有子代数均作为 \(B(H)\) 的具体子集来讨论。

\subsection{弱算子、强算子与范数拓扑}
\label{b2:subsec:D04:01}

\begin{definition}\label{b2:def:operator-topologies}
设 \(H\) 为复 Hilbert 空间，\((T_\alpha)\) 是 \(B(H)\) 中以有向集为指标的网，\(T\in B(H)\)。若对每个 \(\xi\in H\)，都有 \(\bigl\lVert(T_\alpha-T)\xi\bigr\rVert\to0\)，则称 \((T_\alpha)\) 在\term[qiang-suanzi-tuopu]{强算子拓扑}[strong operator topology]中趋于 \(T\)；若对每对 \(\xi,\eta\in H\)，都有 \(\bigl\langle(T_\alpha-T)\xi,\eta\bigr\rangle\to0\)，则称它在\term[ruo-suanzi-tuopu]{弱算子拓扑}[weak operator topology]中趋于 \(T\)。两者简称 SOT、WOT；范数收敛则要求 \(\|T_\alpha-T\|\to0\)。
\end{definition}
\symidx{SOT：强算子拓扑}
\symidx{WOT：弱算子拓扑}

这里的网是以有向集为指标的族，极限要求从某个指标起均满足给定邻域条件；不能一般只用序列描述闭包。SOT 的基本邻域同时限制有限多个 \(\bigl\lVert(S-T)\xi_j\bigr\rVert\)，WOT 则同时限制有限多个矩阵系数。直接由定义和 Cauchy--Schwarz 不等式可得
\[
\text{范数收敛}\ \Longrightarrow\ \text{SOT 收敛}\ \Longrightarrow\ \text{WOT 收敛}.
\]

在 \(\ell^2(\mathbb N)\) 中，令 \(P_n\) 投影到前 \(n\) 个标准基向量。对 \(\xi=(\xi_j)\)，尾和 \(\bigl\lVert(I-P_n)\xi\bigr\rVert^2=\sum_{j>n}|\xi_j|^2\to0\)，故 \(P_n\to I\) 强收敛；但 \(\|I-P_n\|=1\)，并不范数收敛。再令 \(S e_j=e_{j+1}\)。对有限支集向量 \(\xi,\eta\)，\(\langle S^n\xi,\eta\rangle\) 最终为零；利用有限支集向量的稠密性和 \(\|S^n\|=1\)，可推广到任意 \(\xi,\eta\)。因此 \(S^n\to0\) 弱收敛，但 \(\|S^n e_1\|=1\)，不强收敛。

\subsection{双交换子定理}
\label{b2:subsec:D04:02}

设 \(H\) 为复 Hilbert 空间，\(\mathcal S\subseteq B(H)\)。满足与 \(\mathcal S\) 中每个算子交换的全体有界算子组成\term[jiaohuanzi-daishu]{交换子代数}[commutant]
\[
\mathcal S'=\bigl\{\,T\in B(H):TS=ST\text{ 对所有 }S\in\mathcal S\,\bigr\},
\qquad \mathcal S''=(\mathcal S')'.
\]
\symidx{\(\mathcal S'\)：有界算子集合的交换子代数}
\symidx{\(\mathcal S''\)：有界算子集合的双交换子}
这与第十一章的中心化子是同类代数条件，但环境现在是全部有界算子，闭包也有指定的拓扑。

\begin{theorem}[von Neumann 双交换子]\label{b2:thm:operator-bicommutant}
设 \(H\) 为复 Hilbert 空间，\(\mathcal A\subseteq B(H)\) 是含恒等算子 \(I_H\) 的 \(*\)-子代数。记 \(\mathcal A'=\{T\in B(H):Ta=aT\text{ 对所有 }a\in\mathcal A\}\)，\(\mathcal A''=(\mathcal A')'\)。下式中的上标 SOT 和 WOT 分别表示强算子拓扑与弱算子拓扑下的闭包。此时
\[
\mathcal A''=\overline{\mathcal A}^{\mathrm{SOT}}
=\overline{\mathcal A}^{\mathrm{WOT}}.
\]
\end{theorem}
\begin{proof}
固定 \(S\)，等式 \(TS=ST\) 在 WOT 下是闭条件，因为其每个矩阵系数是
\(\langle TS\xi,\eta\rangle-\langle T\xi,S^*\eta\rangle\)，对变量 \(T\) 连续。所以交换子代数为 WOT 闭集。又 \(\mathcal A\subseteq\mathcal A''\)，故
\[
\overline{\mathcal A}^{\mathrm{SOT}}\subseteq
\overline{\mathcal A}^{\mathrm{WOT}}\subseteq\mathcal A''.
\]
只需证明 \(T\in\mathcal A''\) 可以在任意有限组向量上同时由 \(\mathcal A\) 逼近。

给定 \(\xi_1,\ldots,\xi_m\)，在 \(H^m\) 中令
\[
\mathcal M=\overline{\bigl\{\,(a\xi_1,\ldots,a\xi_m):a\in\mathcal A\,\bigr\}},
\]
闭包使用 Hilbert 范数。这是闭线性子空间，对每个对角作用 \(a^{(m)}\) 不变；由于 \(a^*\in\mathcal A\)，它的正交补也不变。故正交投影 \(P\) 与每个 \(a^{(m)}\) 交换。把 \(P\) 按 \(H^m\) 的分量写成算子矩阵 \((P_{ij})\)，交换关系逐条目给出 \(P_{ij}a=aP_{ij}\)，即 \(P_{ij}\in\mathcal A'\)。

因为 \(T\in\mathcal A''\)，它与每个 \(P_{ij}\) 交换，从而 \(T^{(m)}P=PT^{(m)}\)。单位元使 \(\xi=(\xi_1,\ldots,\xi_m)\in\mathcal M\)，所以
\(T^{(m)}\xi=T^{(m)}P\xi=PT^{(m)}\xi\in\mathcal M\)。依照 \(\mathcal M\) 的定义，对任意 \(\varepsilon>0\)，存在同一个 \(a\in\mathcal A\) 使
\[
\sum_{j=1}^m\bigl\lVert(T-a)\xi_j\bigr\rVert^2<\varepsilon^2.
\]
这就使 \(a\) 落在指定 SOT 邻域中，证明 \(T\in\overline{\mathcal A}^{\mathrm{SOT}}\)。
\end{proof}

这称为\term[von-Neumann-shuangjiaohuanzi-dingli]{von Neumann 双交换子定理}[von Neumann bicommutant theorem]。\footnote{冯·诺伊曼（Neumann János Lajos，1903-1957），匈牙利裔美国数学家，奠定算子代数理论，并研究量子力学与计算机结构。}论证中，对伴随封闭用于让正交补不变，单位元用于确保 \(\xi\) 落在循环子空间，有限次放大则用于把“一次逼近一个向量”加强到“同一个算子同时逼近有限多个向量”。可对照 Blackadar\cite[Theorem I.9.1.1; Corollary I.9.1.2]{Blackadar2017OperatorAlgebrasCorrected}，其中还给出非退化作用的版本；此处只使用已经证明的含幺版本。

\begin{definition}\label{b2:def:von-neumann-algebra}
设 \(H\) 为复 Hilbert 空间。若 \(\mathcal M\subseteq B(H)\) 是含恒等算子 \(I_H\) 且在弱算子拓扑下闭的 \(*\)-子代数，则称 \(\mathcal M\) 为\term[von-Neumann-daishu]{von Neumann 代数}[von Neumann algebra]。这里 \(\mathcal M'\) 是 \(\mathcal M\) 在 \(B(H)\) 中的交换子代数，\(\mathcal M''=(\mathcal M')'\)。由\cref{b2:thm:operator-bicommutant}，这等价于 \(\mathcal M=\mathcal M''\)，也等价于含幺 \(*\)-子代数在强算子拓扑下闭。
\end{definition}

\subsection{中心、因子、投影与迹}
\label{b2:subsec:D04:03}

设 \(H\) 为复 Hilbert 空间，\(\mathcal M\subseteq B(H)\) 为 von Neumann 代数。其中心为 \(Z(\mathcal M)=\mathcal M\cap\mathcal M'\)。如果 \(Z(\mathcal M)=\mathbb C I_H\)，则称 \(\mathcal M\) 为\term[yinzi-suanzi]{因子}[factor]。若 \(p\in\mathcal M\) 满足 \(p=p^*=p^2\)，则称 \(p\) 为正交投影；它对应 \(H\) 的一个闭子空间，但并非 \(H\) 的每个闭子空间的投影都属于 \(\mathcal M\)。设 \(\tau\) 是 \(\mathcal M\) 上的态。若对任意 \(a,b\in\mathcal M\) 都有 \(\tau(ab)=\tau(ba)\)，则称 \(\tau\) 为\term[jiantai]{迹态}[tracial state]；若对每个非零 \(a\in\mathcal M\) 还有 \(\tau(a^*a)>0\)，则称它为忠实迹态。

\ProseParagraphBreak
\(B(H)\) 是因子。若 \(T\) 与全部秩一正交投影交换，它保持每条直线，故 \(T\xi=\lambda_\xi\xi\)。对两个线性无关向量再考虑它们的和，得到相同的标量，于是 \(T=\lambda I\)。有限维时 \(M_n(\mathbb C)\) 的归一化迹为忠实迹态；而 \(\bigoplus_{j=1}^rM_{n_j}(\mathbb C)\) 的中心为 \(\mathbb C^r\)，\(r>1\) 时不是因子，其迹态包括
\[
\tau(a_1,\ldots,a_r)=\sum_{j=1}^r t_j\frac{\operatorname{tr}(a_j)}{n_j},
\qquad t_j\ge0,\quad\sum_jt_j=1.
\]
所有 \(t_j>0\) 时该态忠实。无限维时，投影像空间的维数已不能单独刻画它在代数中的大小：比较所用的算子是否属于当前代数，同样是条件的一部分。

\subsection{投影等价与因子的类型}
\label{b2:subsec:D04:04}

一个闭子空间可以与自己的真子空间等距同构，例如单边移位把 \(\ell^2(\mathbb N)\) 等距映入去掉第一个坐标的子空间。投影的有限性要问的是：这种等距映射能否由当前 von Neumann 代数中的算子实现。

\begin{definition}\label{b2:def:projection-equivalence}
设 \(H\) 为复 Hilbert 空间，\(\mathcal M\subseteq B(H)\) 为 von Neumann 代数。若 \(v\in\mathcal M\) 满足 \(vv^*v=v\)，则称 \(v\) 为\term[bufen-dengju]{部分等距算子}[partial isometry]。此时 \(v^*v\) 与 \(vv^*\) 都是投影，分别称为它的初投影与终投影。对投影 \(p,q\in\mathcal M\)，若存在部分等距算子 \(v\in\mathcal M\)，使
\[
v^*v=p,\qquad vv^*=q,
\]
则称 \(p\) 与 \(q\) 在 \(\mathcal M\) 中等价，记为 \(p\sim_{\mathcal M}q\)。这一关系称为\term[Murray-von-Neumann-dengjia]{Murray--von Neumann 投影等价}[Murray--von Neumann equivalence]。\footnote{默里（Francis Joseph Murray，1911-1996），美国数学家，与 von Neumann 合作奠定算子环及因子理论。}
\end{definition}

由 \(vv^*v=v\) 可直接验证两个自伴算子的平方仍为自身。算子 \(v\) 在 \(pH\) 上保持范数，像恰为 \(qH\)，在 \((1-p)H\) 上为零。取 \(v=p\) 给出自反性，取伴随给出对称性；若 \(v\) 把 \(p\) 连接到 \(q\)，而 \(w\) 把 \(q\) 连接到 \(r\)，则 \((wv)^*(wv)=p\)、\((wv)(wv)^*=r\)，得到传递性。下文写 \(p\sim q\) 时，等价都在指定的 \(\mathcal M\) 内。

\begin{definition}\label{b2:def:finite-minimal-projection}
设 \(H\) 为复 Hilbert 空间，\(\mathcal M\subseteq B(H)\) 为 von Neumann 代数，\(p\in\mathcal M\) 为投影。对另一投影 \(q\in\mathcal M\)，规定 \(q\le p\) 表示 \(pq=qp=q\)，即 \(qH\subseteq pH\)。若对每个满足 \(q\le p\) 且 \(q\sim_{\mathcal M}p\) 的投影 \(q\)，都有 \(q=p\)，则称 \(p\) 为\term[youxian-touying]{有限投影}[finite projection]；否则称它为\term[wuxian-touying]{无限投影}[infinite projection]。若 \(p\ne0\)，且它的子投影只有 \(0\) 和 \(p\)，则称 \(p\) 为\term[jixiao-touying]{极小投影}[minimal projection]。
\end{definition}

极小投影必为有限投影，因为与非零投影等价的投影仍非零。“有限”取决于环境代数；例如在无限维 \(H\) 上，\(I_H\) 在 \(\mathbb C I_H\) 内为极小投影，却在 \(B(H)\) 内为无限投影。

\begin{proposition}\label{b2:prop:finite-projection-models}
设 \(H\ne0\) 为复 Hilbert 空间。\(B(H)\) 中的投影有限，当且仅当它的像空间有限维；非零投影极小，当且仅当它的像空间一维。若 von Neumann 代数 \(\mathcal M\subseteq B(H)\) 有忠实迹态，则它的每个投影均有限。
\end{proposition}
\begin{proof}
在 \(B(H)\) 中，\(p\sim q\) 正好表示 \(pH\) 与 \(qH\) 酉同构：一个方向由部分等距算子的限制给出，另一个方向把酉同构在 \((1-p)H\) 上延拓为零。有限维空间不能与自己的真子空间酉同构，所以有限秩投影有限。

若 \(pH\) 无限维，在其中取正交单位向量 \(\xi_1,\xi_2,\ldots\)。令 \(v\xi_j=\xi_{j+1}\)，在这列向量闭线性张成空间于 \(pH\) 内的正交补上令 \(v\) 为恒等，在 \((1-p)H\) 上令它为零。于是 \(v^*v=p\)，而 \(vv^*=p-e\)，其中 \(e\) 为到 \(\mathbb C\xi_1\) 的投影。故 \(p\) 为无限投影。一维空间没有非零真子空间，维数大于一时则有，所以极小性也得到刻画。

最后设 \(\tau\) 为 \(\mathcal M\) 上的忠实迹态，\(q\le p\) 且 \(v^*v=p\)、\(vv^*=q\)。迹性质给出 \(\tau(p)=\tau(q)\)，从而 \(\tau(p-q)=0\)。由于 \(p-q\) 为投影，忠实性迫使 \(p-q=0\)。
\end{proof}

\begin{definition}\label{b2:def:factor-types}
设 \(H\ne0\) 为复 Hilbert 空间，\(\mathcal M\subseteq B(H)\) 为 von Neumann 代数，并且 \(Z(\mathcal M)=\mathbb C I_H\)，即 \(\mathcal M\) 为因子。按照 \(\mathcal M\) 中投影的性质，定义下列类型。
\begin{center}
\begin{tabularx}{\linewidth}{@{}lX@{}}
\toprule
类型 & 投影条件 \\
\midrule
\(\mathrm I\) 型 & 存在非零极小投影。 \\
\(\mathrm{II}\) 型 & 不存在非零极小投影，但存在非零有限投影。 \\
\(\mathrm{III}\) 型 & 不存在非零有限投影。 \\
\bottomrule
\end{tabularx}
\end{center}
对于 \(\mathrm{II}\) 型因子，若单位投影有限，则称为 \(\mathrm{II}_1\) 型；若单位投影无限，则称为 \(\mathrm{II}_\infty\) 型。
\end{definition}

三种条件互相排斥并穷尽所有因子。它们区分的是投影结构，并未给出同一类型内部的同构分类。这里采用因子的投影判据；一般 von Neumann 代数的类型通过中心投影分解定义，各个非零中心分量可以属于不同类型。投影等价、有限性和一般的中心分解见 Blackadar\cite[\nopp III.1.1.2; Definition III.1.3.1; III.1.4.5--III.1.4.7]{Blackadar2017OperatorAlgebrasCorrected}。

\ProseParagraphBreak
由\cref{b2:prop:finite-projection-models}，\(B(H)\) 总是 \(\mathrm I\) 型因子。\(M_n(\mathbb C)\) 属于 \(\mathrm I_n\) 型，而 \(B(\ell^2(\mathbb N))\) 属于 \(\mathrm I_\infty\) 型；前者的单位投影有限，后者的单位投影无限。一般 \(\mathrm I\) 型因子也与某个 \(B(K)\) 同构，其结构证明是把等价的极小投影组织成矩阵单位，见\cite[\nopp III.1.5.2--III.1.5.3; III.1.5.12]{Blackadar2017OperatorAlgebrasCorrected}。当 \(K\) 不可分时，还可用其 Hilbert 维数作为 \(\mathrm I\) 型的下标。

\ProseParagraphBreak
\(\mathrm{II}_1\) 型也有有限单位投影，却没有矩阵代数中的极小投影；有限投影可以继续分解，不能把它理解成有限秩投影。\(\mathrm{II}_\infty\) 型的单位投影无限，但代数内部仍有非零有限投影；\(\mathrm{III}\) 型连这样的有限部分也不存在。尤其不能仅由“没有迹态”断定一个因子为 \(\mathrm{III}\) 型：\(B(\ell^2(\mathbb N))\) 就没有迹态。若它有迹态 \(\tau\)，取 \(v_1e_j=e_{2j}\)、\(v_2e_j=e_{2j-1}\)，则 \(v_i^*v_i=I\)，两个终投影之和为 \(I\)，因而 \(1=\tau(I)=\tau(v_1v_1^*)+\tau(v_2v_2^*)=2\)，矛盾。\(\mathrm{II}\) 与 \(\mathrm{III}\) 型的具体构造、非有界迹以及 \(\mathrm{III}\) 型的进一步细分见\cite[\nopp III.2.5; III.3; III.4.6]{Blackadar2017OperatorAlgebrasCorrected}。
